% Requires \usepackage{booktabs} and \usepackage{tabularx}.
\begin{table}[ht]
  \centering
  \caption{Mechanism-level comparison of the four trackers.}
  \label{tab:tracker_arch}
  \small
  \setlength{\tabcolsep}{4pt}
  \renewcommand{\arraystretch}{1.25}
  \begin{tabularx}{\columnwidth}{l *{4}{>{\centering\arraybackslash}X}}
    \toprule
     & ByteTrack & BoT-SORT & OC-SORT & Deep OC-SORT \\
    \midrule
    Year / lineage
     & 2022 / SORT
     & 2022 / ByteTrack
     & 2023 / SORT (OC)
     & 2023 / OC-SORT \\
    \addlinespace
    Motion model
     & Linear KF
     & Improved KF (w,h)
     & Obs.-centric KF
     & Obs.-centric KF \\
    \addlinespace
    Appearance / ReID
     & None
     & VeRi-776 CLIP-ReID
     & None
     & VeRi-776 CLIP-ReID \\
    \addlinespace
    Camera-motion comp.
     & No
     & Yes (GMC)
     & No
     & Yes (CMC) \\
    \addlinespace
    Association
     & IoU
     & IoU + appearance
     & IoU + OC momentum
     & Motion + adapt. appearance \\
    \addlinespace
    Low-score (BYTE)
     & Yes
     & Yes
     & No
     & No \\
    \bottomrule
  \end{tabularx}
  \par\smallskip
  \footnotesize All four trackers share a single set of RT-DETR-L detections (``detect once, track four times'') ingested at \texttt{min\_conf}=0.1, so reported differences are purely in association. The roster forms a 2$\times$2 design over motion model (linear vs.\ observation-centric Kalman filter) and appearance (motion-only vs.\ ReID-augmented); ReID, where used, is the VeRi-776 CLIP-ReID vehicle encoder (FP16). ``Obs.-centric KF'' denotes OC-SORT's observation-centric re-update, momentum, and recovery (ORU/OCM/OCR). ``Low-score (BYTE)'' marks the two-stage association that also matches low-confidence detections (\texttt{use\_byte}; disabled for the OC-SORT family in our configs).
\end{table}
